% The front matter as data: title, authors, their e-mail, ORCID and contributions, and the one affiliation.
% No journal's syntax appears here, because every class asks for the same facts in a different shape: one wants
% a single name per author with the address as key=value pairs, another wants the given and family names split
% into separate commands and the address broken into fields of its own, and the reading copy wants one
% \and-separated line. So this file states each fact once and a template renders it, exactly as the abstract and
% the keywords already work. Supporting one more journal means writing the renderers below in the new
% template, and changing nothing here.
%
% A template must define, before it uses \projauthorlist:
%   \projauthor{initials}{given}{family}{email}{orcid}{CRediT roles}      one non-corresponding author
%   \projauthorcorr{...}                                                 same six, for the corresponding author
% ORCID may be empty, so a renderer that prints it guards the field with \projifhasorcid.
%
% THIS IS THE FIRST FILE TO EDIT IN A NEW PAPER. Nothing else holds any of these facts.

\newcommand{\projtitle}{The Title of the Paper, Stated in Full and Without a Subtitle Unless One Earns Its Place}
\newcommand{\projshorttitle}{Short Running Title}
\newcommand{\projrunningauthors}{Family-Name et al., 2026}

\newcommand{\projaffildiv}{Department or Institute}
\newcommand{\projaffilorg}{University}
\newcommand{\projaffilstreet}{Street and Number}
\newcommand{\projaffilcity}{City}
\newcommand{\projaffilstate}{Region}
\newcommand{\projaffilpostcode}{00000}
\newcommand{\projaffilcountry}{Country}

% Not every author has an ORCID on file, and a renderer that prints the field unconditionally would emit an empty
% link for the ones who do not. This runs its second argument only when the first is a non-empty ORCID.
\newcommand{\projifhasorcid}[2]{\def\projorcidarg{#1}\ifx\projorcidarg\empty\else#2\fi}

% The authors in the order they appear on the paper. The corresponding author is marked by which macro the line
% calls, not by a flag, so no layout has to decide what "corresponding" looks like in its own notation.
% The roles are the CRediT taxonomy's own terms: they go into the author block of a class that has one, and into
% the contributions declaration, which is rendered from this same list.
\newcommand{\projauthorlist}{%
	\projauthorcorr{A.B.}{Ada}{Firstauthor}{ada.firstauthor@example.org}{0000-0000-0000-0000}%
		{Conceptualization, Methodology, Software, Formal analysis, Writing -- original draft}%
	\projauthor{C.D.}{Carl}{Secondauthor}{carl.secondauthor@example.org}{}%
		{Validation, Writing -- review \& editing, Supervision}%
}
